\documentclass[11pt,letterpaper]{article}
\usepackage[T1]{fontenc}
\usepackage[utf8]{inputenc}
\usepackage{lmodern}
\usepackage[margin=1in]{geometry}
\usepackage{booktabs}
\usepackage{longtable}
\usepackage{enumitem}
\usepackage{listings}
\usepackage{xcolor}
\usepackage[hidelinks]{hyperref}

\definecolor{backcolour}{rgb}{0.96,0.96,0.94}
\definecolor{cautionred}{rgb}{0.55,0,0}

% Listings are configured so that anything inside one can be selected from the
% PDF and pasted straight into a terminal:
%   numbers=none        no line numbers to catch in the selection
%   columns=fullflexible no padding spaces inserted between characters
%   upquote=true        straight quotes, not typographic ones
%   breaklines=false    a wrapped line would paste with a spurious newline,
%                       so every command below is kept short enough to fit
\lstset{
    basicstyle=\ttfamily\small,
    backgroundcolor=\color{backcolour},
    numbers=none,
    frame=single,
    framesep=5pt,
    columns=fullflexible,
    keepspaces=true,
    upquote=true,
    showstringspaces=false,
    breaklines=false,
    xleftmargin=0pt,
    aboveskip=8pt,
    belowskip=8pt
}

\newenvironment{caution}[1]{%
    \par\addvspace{8pt}%
    \noindent{\color{cautionred}\rule{\linewidth}{0.8pt}}\par\vspace{2pt}
    \noindent\textbf{\color{cautionred}Caution: #1}\par\vspace{2pt}
}{%
    \par\vspace{2pt}\noindent{\color{cautionred}\rule{\linewidth}{0.4pt}}%
    \par\addvspace{8pt}%
}

\title{BASIC CAN COMMANDS \\ \large Arctos Robotic Arm --- terminal quick reference}
\date{September 2026}

\begin{document}
\maketitle

Everything here is a command you can select from this PDF and paste directly into
a terminal. Nothing needs editing except the CAN ID, and Section~\ref{sec:frames}
gives the ready-made frame for every ID on this arm so you never have to work out
a checksum by hand.

For why any of this works, see \texttt{Arctos\_CAN\_Interface.pdf}. For bringing up
a specific joint from scratch, see the two SOP documents.

\tableofcontents
\newpage

\section{Set the interface up}

The CANable presents as a serial device. \texttt{slcand} bridges it to a Linux
network interface, and this has to be redone every time the adapter is plugged in
or re-enumerates. Run these in order.

\subsection{Check the adapter is there}
\begin{lstlisting}
lsusb
\end{lstlisting}

Expect a line containing \texttt{16d0:117e}. Then find the stable device name:

\begin{lstlisting}
ls -l /dev/serial/by-id/
\end{lstlisting}

Use the \texttt{/dev/serial/by-id/} name rather than \texttt{/dev/ttyACM0}. The
\texttt{ttyACM} number changes whenever the adapter re-enumerates; the by-id name
does not.

\subsection{Bring it up at 500 kbit/s}
\begin{lstlisting}
sudo pkill slcand
sudo ip link delete can0 2>/dev/null
sudo slcand -o -c -s6 /dev/serial/by-id/usb-Openlight_Labs_CANable2* can0
sudo ip link set can0 txqueuelen 1000
sudo ip link set up can0
\end{lstlisting}

\begin{caution}{The \texttt{-s6} flag is a coded index, not a number}
\texttt{-s0} is 10k, \texttt{-s1} 20k, \texttt{-s2} 50k, \texttt{-s3} 100k,
\texttt{-s4} 125k, \texttt{-s5} 250k, \textbf{\texttt{-s6} 500k}, \texttt{-s7}
800k, \texttt{-s8} 1M. This arm runs at 500 kbit/s so \texttt{-s6} is correct.
Older scripts in this project used \texttt{-s3} and silently attached at
100 kbit/s, which looks like a dead bus.
\end{caution}

\subsection{Verify}
\begin{lstlisting}
ip -br link show can0
\end{lstlisting}

Expect \texttt{can0} followed by \texttt{UP}. That is the whole check.

\begin{caution}{Do not try to set or read the bitrate with \texttt{ip}}
This does not work on this hardware and will fail:
\begin{lstlisting}
sudo ip link set can0 up type can bitrate 500000
\end{lstlisting}
That netlink attribute is for native CAN controllers, not slcan. For the same
reason \texttt{ip -details link show can0} reports \texttt{bitrate 0} and
\texttt{clock 0} on a perfectly healthy adapter. That is normal, and it is not a
fault worth chasing. The bitrate was fixed at attach time by \texttt{-s6}.
\end{caution}

\section{Shut it down so you can unplug}

In this order. Stop any running control script first with Ctrl+C.

\begin{lstlisting}
sudo ip link set can0 down
sudo pkill slcand
\end{lstlisting}

Then unplug the adapter. If a motor is energized and it is safe to release it,
disable the coils first --- see Section~\ref{sec:frames}.

\begin{caution}{Dropping coils on a gravity-loaded joint lets it fall}
Disabling a motor removes holding torque. On a bench motor that is harmless; on an
assembled arm the joint may drop. Support it first, or leave it holding.
\end{caution}

\section{Ready-made frames for every joint}
\label{sec:frames}

Every frame is \texttt{[command][data...][checksum]}, and the checksum includes the
CAN ID:

\begin{lstlisting}
checksum = (can_id + sum_of_all_other_bytes) & 0xFF
\end{lstlisting}

Because the ID is in the sum, \textbf{a frame copied from another joint is
silently rejected}. The table below has the arithmetic already done. Find your
joint's ID and use that row.

\begin{center}
\begin{tabular}{lllllll}
\toprule
\textbf{ID} & \textbf{Read enc.} & \textbf{Enable} & \textbf{Disable}
            & \textbf{E-stop} & \textbf{Enable state} & \textbf{Error} \\ \midrule
\texttt{0x01} & \texttt{001\#3132} & \texttt{001\#F301F5} & \texttt{001\#F300F4} & \texttt{001\#F7F8} & \texttt{001\#3A3B} & \texttt{001\#393A} \\
\texttt{0x02} & \texttt{002\#3133} & \texttt{002\#F301F6} & \texttt{002\#F300F5} & \texttt{002\#F7F9} & \texttt{002\#3A3C} & \texttt{002\#393B} \\
\texttt{0x03} & \texttt{003\#3134} & \texttt{003\#F301F7} & \texttt{003\#F300F6} & \texttt{003\#F7FA} & \texttt{003\#3A3D} & \texttt{003\#393C} \\
\texttt{0x04} & \texttt{004\#3135} & \texttt{004\#F301F8} & \texttt{004\#F300F7} & \texttt{004\#F7FB} & \texttt{004\#3A3E} & \texttt{004\#393D} \\
\texttt{0x05} & \texttt{005\#3136} & \texttt{005\#F301F9} & \texttt{005\#F300F8} & \texttt{005\#F7FC} & \texttt{005\#3A3F} & \texttt{005\#393E} \\
\texttt{0x06} & \texttt{006\#3137} & \texttt{006\#F301FA} & \texttt{006\#F300F9} & \texttt{006\#F7FD} & \texttt{006\#3A40} & \texttt{006\#393F} \\
\bottomrule
\end{tabular}
\end{center}

\subsection{Emergency stop}

Work out your joint's e-stop line from the table and keep it pasted in a second
terminal \emph{before} you command any motion. For the joint most recently worked
on, CAN ID \texttt{0x06}:

\begin{lstlisting}
cansend can0 006#F7FD
\end{lstlisting}

\begin{caution}{Check the e-stop is addressed to the right joint}
An earlier version of this project's documentation gave a joint's emergency stop
as \texttt{cansend can0 001\#F7F8} when the motor was not on ID \texttt{0x01}. It
would have done nothing at all. Confirm the ID by scan, then take the matching row
above.
\end{caution}

\subsection{Enable and disable coils}

For CAN ID \texttt{0x06}:

\begin{lstlisting}
cansend can0 006#F301FA
cansend can0 006#F300F9
\end{lstlisting}

The first enables, the second disables. An enabled joint holds position under power
and can move; keep clear of the mechanism before enabling.

\section{Watch what is on the bus}

Everything:

\begin{lstlisting}
candump can0
\end{lstlisting}

One joint only, by ID. This example filters to \texttt{0x06}:

\begin{lstlisting}
candump can0,006:7FF
\end{lstlisting}

Log to a file while you work:

\begin{lstlisting}
candump -L can0 > can_log.txt
\end{lstlisting}

Run \texttt{candump} in a second terminal while sending commands in the first. It
is the fastest way to see whether a frame went out and whether anything answered.

\section{Read-only queries}

None of these move the motor or change its state, so they are safe to send at any
time. Replace \texttt{006} and the checksum with your joint's row from
Section~\ref{sec:frames}.

\begin{center}
\begin{tabular}{lll}
\toprule
\textbf{What} & \textbf{Opcode} & \textbf{Expected at rest} \\ \midrule
Encoder position   & \texttt{0x31} & stable, no jitter, no drift \\
Velocity           & \texttt{0x32} & zero \\
Position error     & \texttt{0x39} & small, tens of counts \\
Enable state       & \texttt{0x3A} & \texttt{01} enabled, \texttt{00} disabled \\
Shaft protection   & \texttt{0x3E} & \texttt{00} \\
Liveness           & \texttt{0xF1} & answers \texttt{F1 01} \\
\bottomrule
\end{tabular}
\end{center}

Read the encoder on \texttt{0x06}:

\begin{lstlisting}
cansend can0 006#3137
\end{lstlisting}

The reply is \texttt{[0x31][6 bytes signed position][checksum]}, big-endian, on the
16384-counts-per-motor-revolution scale. To turn it into joint degrees:

\begin{lstlisting}
joint_degrees = raw48 * 360.0 / 16384 / gear_ratio
\end{lstlisting}

\section{Scan for what is actually on the bus}

Never trust a config file or a table for a CAN ID. This sends only the read-only
\texttt{0x31} query, never enables coils, and never commands motion:

\begin{lstlisting}
cd ~/arctos_ws/src/arctos_can_control/arctos_can_control
python3 can_id_scan.py --channel can0 --ids 1-16
\end{lstlisting}

Exactly one responder is the healthy result. Two or more means another driver is
powered on the bus with a colliding ID. None, with the board powered and wiring
confirmed, is the signature of an RS485 board that has no CAN transceiver at all.

\section{A safe test sequence}

In order, from cold. Stop at the first step that does not do what it should.

\begin{enumerate}[leftmargin=*]
    \item Bring the interface up, confirm \texttt{can0 UP}.
    \item Start \texttt{candump can0} in a second terminal.
    \item Scan for the real CAN ID.
    \item Read the encoder twice and confirm the value is stable.
    \item Read the enable state and the shaft-protection state.
    \item Paste your joint's e-stop line into a third terminal, ready but unsent.
    \item Only now enable the coils, and only if the mechanism is clear.
    \item Command one small move and confirm the encoder changed by what you asked.
    \item Disable the coils.
\end{enumerate}

\section{When something is wrong}

\begin{longtable}{p{0.36\textwidth}p{0.58\textwidth}}
\toprule
\textbf{Symptom} & \textbf{What to do} \\ \midrule
\endhead
\texttt{can0} does not exist &
  \texttt{slcand} died, usually because the adapter re-enumerated. Re-run the
  bring-up. Check \texttt{/dev/serial/by-id/} for the current name. \\
Adapter is plugged in but \texttt{can0} is missing &
  Expected. \texttt{slcand} does not restart itself after a replug. \\
\texttt{ip} shows \texttt{bitrate 0} &
  Normal on slcan. Not a fault. \\
Nothing answers any ID &
  Check panel CAN ID, CAN rate and CANRSP; check CAN\_H, CAN\_L and termination;
  confirm the driver has power. If still silent, check whether the board is the
  RS485 variant. \\
The command was accepted but nothing moved &
  The bus cannot tell ``not driven'' from ``driven but cannot move''. Disable the
  coils and turn the joint by hand. \\
A move under-delivers, then delivers nothing &
  Stop. Do not resend --- that only heats the motor. Back off and re-attempt once,
  or investigate the mechanism. \\
Motor is hot &
  Stop and let it cool. There is no temperature-read opcode, so nothing in software
  will warn you. Warm is normal; too hot to touch is not. \\
\bottomrule
\end{longtable}

\end{document}
